\documentclass[12pt]{article}
\usepackage[utf8]{inputenc}
\usepackage[T1]{fontenc}
\usepackage{hyperref}
\title{Sindhi Bhil Language Resources for K-12 Curriculum Development}
\author{Compiled by Claude Sonnet 5.5 (claude-sonnet-5-5) for Andr\'as Kornai}
\date{}
\begin{document}
\maketitle

\section{Introduction}
Sindhi Bhil (Sindhi Bhili, Bhilki; ISO 639-3 sbn, Glottocode sind1270) is an Indo-Aryan variety closely tied to Sindhi, spoken by Hindu communities, mainly Bhils and Meghwars, in Sindh, with some speakers in Balochistan and South Punjab in Pakistan. Wikipedia treats it as a dialect of Sindhi. Speaker estimates range from about 56,000 to 86,500. Listed dialects are Badin, Mohrano, Nuclear Sindhi Bhil and Sindhi Meghwar. It is written in the Arabic script, with Devanagari used mainly for religious purposes. Search found only evangelistic audio, with no school materials.

\section{Primary Educational Resources}
No primer, textbook or school programme specific to Sindhi Bhil turned up. Standard Sindhi is the closest school language, and its materials were not searched here.

\section{Linguistic and Reference Materials}
Only encyclopedic descriptions turned up. No Sindhi Bhil dictionary or grammar was found.
\begin{itemize}
\item Wikipedia, Sindhi Bhili. \url{https://en.wikipedia.org/wiki/Sindhi_Bhili}
\item Wikipedia, Sindhi Bhils. \url{https://en.wikipedia.org/wiki/Sindhi_Bhils}
\item Wikipedia, Bhil. \url{https://en.wikipedia.org/wiki/Bhil}
\end{itemize}

\section{Religious and Community Resources}
No Sindhi Bhil scripture translation was found. Joshua Project has a profile of the Bhil in Pakistan and of the related Megh people. The Sindhi Bible resources it points to were not confirmed as being in the Bhil variety.
\begin{itemize}
\item Joshua Project, Bhil in Pakistan. \url{https://joshuaproject.net/index.php/people_groups/16414/PK}
\item Joshua Project, Megh (Muslim traditions) in Pakistan. \url{https://joshuaproject.net/people_groups/17619/PK}
\end{itemize}

\section{Digital and Audio Resources}
Global Recordings Network has a Sindhi Bhil language page and a ``Words of Life'' programme of about an hour, plus a Sindhi Meghwar variety page with a 40-part illustrated Bible lesson series. These are evangelistic in purpose, but they are the only recorded native-language material found. Downloads are offered as MP3 and slideshow.
\begin{itemize}
\item Global Recordings Network, Sindhi Bhil. \url{https://globalrecordings.net/en/language/sbn}
\item Global Recordings Network, Sindhi Meghwar. \url{https://globalrecordings.net/en/language/16736}
\item Global Recordings Network, Words of Life with Sindhi, Hindu and Hindi. \url{https://globalrecordings.net/en/program/10230}
\item Global Recordings Network, Sindhi: Shikari Bhil. \url{https://www.globalrecordings.net/en/language/26102}
\end{itemize}

\section{Curriculum Development Context}
The language is not known to be taught in schools. Because it is so close to Sindhi, a curriculum could use standard Sindhi materials with Bhil-specific vocabulary and oral texts added. Mohrano may be a separate language, and the dialects differ, so a variety must be chosen first. The audio recordings could supply oral-text material, subject to checking the licence and the content. This dossier is short because little was found.
\end{document}